%% ============================================================
%% MODELO DE ATIVIDADE · ATIVIDADE COM USUÁRIOS E EQUIPE DE SAÚDE
%%
%% Copie este arquivo para atividades/ (atividade-02.tex, por exemplo),
%% inclua-o no relatorio-aacc.tex com \input{atividades/atividade-02}
%% e escreva. Cada \preencher{…} é uma pendência: troque-o pelo seu
%% texto. Cada \figuraprovisoria também: troque-a pela figura de verdade.
%%
%% Uma atividade é UMA entrega sua, com começo, meio e fim. "Participei
%% do projeto X por dois anos" não é uma atividade: são várias. Divida.
%% ============================================================
\begin{atividade}{Sessões com os usuários do sistema S}{
  tipo         = clinico,
  modalidade   = {},   % a chave da modalidade pedida: EXT, PES, COMP, VPC, EVT, PRD, ORG, ENS, GES, CUR, OUT
  alternativa  = {},   % a modalidade alternativa, se o Colegiado não aceitar a primeira
  horas        = {},   % as horas desta atividade, com base nos registros (Parte IV)
  inicio       = {},   % AAAA-MM
  fim          = {},   % AAAA-MM, ou: em andamento
  projeto      = {},
  papel        = {},   % Responsável, Corresponsável ou Colaborador(a)
  supervisor   = {},   % quem confirma: nome e cargo (o supervisor do projeto, o gerente)
  registros    = {},   % os códigos no SOMA, com ponto e vírgula: NBL-14; NRO-PRO-003-7
  competencias = {},   % as competências das DCN que ela desenvolveu (I a VIII): I, III, V
  disciplinas  = {},   % as disciplinas do seu curso ligadas a ela: ELE075 Eletrônica; ELE067 Circuitos
}

\begin{contexto}
\preencher{Quem são os usuários, do que precisavam e como a equipe chegou até eles.}
\end{contexto}

\begin{contribuicao}
\preencher{O seu papel técnico nas sessões: configuração, ajuste, coleta de dados, segurança.}
\end{contribuicao}

\begin{desenvolvimento}
\fichatecnica{
  local      = {},   % Local
  publico    = {},   % Público atendido
  sessoes    = {},   % Sessões (quantidade e duração)
  supervisao = {},   % Supervisão profissional
  etica      = {},   % Aprovação ética e consentimento
}

\preencher{O protocolo das sessões, a supervisão profissional, os cuidados de segurança e com os dados (LGPD). Nada que identifique uma pessoa.}
\end{desenvolvimento}

\begin{resultados}
\preencher{O que as sessões mostraram --- sem identificar ninguém --- e o que mudou no projeto por causa delas.}
\end{resultados}

\begin{evidencias}
% Uma linha por evidência: \evidencia{tipo}{identificador}{o que mostra}{onde conferir}
% \evidencia{Registro}{<projeto>/sessoes/<AAAA-MM>}{O registro das sessões (sem dados pessoais)}{SOMA / Depto. Clínico (a pedido)}
\end{evidencias}

\begin{aprendizados}
\preencher{O que o contato com o usuário real ensinou sobre projeto de engenharia.}
\end{aprendizados}
\end{atividade}
